\section*{Chapitre \ref{ch:g5:weather-and-climate} --- Le temps et le climat}

\begin{solution}{exo:g5:weather-and-climate:1}
(a) Temps (ce matin seulement). (b) \omterm{def:g5:weather-and-climate:climate}{Climat} (chaque été). (c) Temps (hier
après-midi). (d) \omterm{def:g5:weather-and-climate:climate}{Climat} (chaque mois, année après année).
\end{solution}

\begin{solution}{exo:g5:weather-and-climate:2}
Le \omterm{def:g5:weather-and-climate:climate}{climat} d'un lieu est le temps qu'il y fait habituellement, mois après
mois, sur de nombreuses années. Les météorologues font la moyenne des
mesures de 30 ans.
\end{solution}

\begin{solution}{exo:g5:weather-and-climate:3}
Le mois le plus froid est décembre (environ 6~$^\circ$C~; janvier est
presque aussi froid). Le mois le plus humide est novembre (environ
160~mm). Non~: la courbe ne descend jamais sous 0~$^\circ$C.
\end{solution}

\begin{solution}{exo:g5:weather-and-climate:4}
Juillet est le mois le plus chaud, avec environ 5~$^\circ$C. Aucun mois ne
dépasse 10~$^\circ$C.
\end{solution}

\begin{solution}{exo:g5:weather-and-climate:5}
L'\omterm{def:g3:continents-and-oceans:equator}{équateur}~: la zone chaude. Le \omterm{def:g3:continents-and-oceans:poles}{pôle Nord}~: la zone froide (celle du nord).
À mi-chemin entre l'\omterm{def:g3:continents-and-oceans:equator}{équateur} et le \omterm{def:g3:continents-and-oceans:poles}{pôle Sud}~: la zone tempérée du sud
(elle s'étend entre le \omterm{def:g5:weather-and-climate:tropics}{tropique} du Capricorne et le \omterm{def:g5:weather-and-climate:tropics}{cercle polaire}
antarctique).
\end{solution}

\begin{solution}{exo:g5:weather-and-climate:6}
Hilo~: la zone chaude. Phoenix et Seattle~: la zone tempérée du nord.
Utqiaġvik~: la zone froide du nord, au-delà du \omterm{def:g5:weather-and-climate:tropics}{cercle polaire} arctique.
\end{solution}

\begin{solution}{exo:g5:weather-and-climate:7}
Année entière~: $147 + 96 + 106 + 81 + 48 + 37 + 15 + 25 + 41 + 99 + 160 + 145
= 1\,000$ millimètres. Décembre, janvier et février~:
$145 + 147 + 96 = 388$ millimètres, plus du tiers de la \omterm{def:g4:clouds-rain-snow:precipitation}{pluie} de l'année en
trois mois.
\end{solution}

\begin{solution}{exo:g5:weather-and-climate:8}
Phoenix~: $35 - 14 = 21$ degrés. Hilo~: $25 - 22 = 3$ degrés. Les saisons
font beaucoup plus varier la température à Phoenix~; à Hilo, elle est
presque la même toute l'année.
\end{solution}

\begin{solution}{exo:g5:weather-and-climate:9}
Un \omterm{def:g5:weather-and-climate:climate}{climat} de \omterm{def:g5:weather-and-climate:rainforest}{forêt tropicale humide}~: il fait chaud toute l'année (27 ou
28~$^\circ$C, sans hiver), et il tombe beaucoup de \omterm{def:g4:clouds-rain-snow:precipitation}{pluie} chaque mois (au
moins 160~mm, et plus de 2\,000~mm sur l'année). La ville se trouve très
probablement dans la zone chaude, près de l'\omterm{def:g3:continents-and-oceans:equator}{équateur}.
\end{solution}

\begin{solution}{exo:g5:weather-and-climate:10}
Première explication~: la plage est dans la zone chaude, où il fait chaud
toute l'année, et le jardin enneigé dans une zone tempérée ou froide, où
janvier tombe en hiver. Seconde explication~: la plage est dans la moitié
sud de la Terre, où janvier tombe en été, et le jardin enneigé dans la
moitié nord, où janvier tombe en hiver~: les saisons sont inversées entre
les deux moitiés du \omterm{def:g3:continents-and-oceans:globe}{globe}.
\end{solution}

\begin{solution}{exo:g5:weather-and-climate:11}
Températures~: Phoenix est beaucoup plus chaud, environ 35~$^\circ$C en
juillet contre environ 20~$^\circ$C en août à Seattle, et son mois de
janvier (environ 14~$^\circ$C) est déjà aussi chaud que le mois de mai de
Seattle. \omterm{def:g4:clouds-rain-snow:precipitation}{Pluie}~: moins de 200~mm par an à Phoenix, un \omterm{def:g5:weather-and-climate:desert}{désert}, contre
environ 1\,000~mm à Seattle. Elles n'ont donc pas le même \omterm{def:g5:weather-and-climate:climate}{climat}. Les zones
nous disent, en gros, à quel point le Soleil chauffe un lieu, de l'\omterm{def:g3:continents-and-oceans:equator}{équateur}
aux \omterm{def:g3:continents-and-oceans:poles}{pôles}~; elles ne nous disent ni combien il pleut, ni l'effet de la \omterm{def:g2:water-on-earth:ocean}{mer},
des \omterm{def:g1:landscapes-around-us:landform}{montagnes} ou des vents.
\end{solution}
